\subsection{Sumary about the Theoretical Background}
The theoretical frameworks and technological paradigms examined in this chapter provide the essential blueprint for aBridgeAI. By grounding the system in the principles of CLT and individualized learner modeling, the system establishes a cognitively aware baseline that respects human memory constraints. The integration of a hybrid SM-2 algorithm and dynamic evaluation rubrics translates these psychological principles into an actionable, mathematical engine capable of adaptive spacing. Furthermore, LLMs, RAG pipelines, and Vector Databases supports the requisite technological infrastructure to maintain semantic context and relational depth at scale. 

Together, these theoretical and computational elements form a cohesive ecosystem that directly addresses the limitations of traditional, isolated SRS.

Having established why aBridgeAI is designed the way it is and the core technologies that power it, the focus must now shift to implementation. Chapter 4 – System Requirements and Analysis will translate these theoretical foundations into concrete software requirements, architectural diagrams, and functional specifications, detailing exactly how this cognitive ecosystem is operationalized in practice.